\chapter{Stochastic Volatility}\label{ch:dv:stochastic-volatility}

Asked why the calibrated correlation of her index model is $-0.7$, an options trader answers
without a formula: because when the index falls, volatility rises. Asked why the \omterm{def:dv:stochastic-volatility:sv}{volatility of
volatility} is 0.6, she points at the wings of the one-year smile, which the model can only
reach if variance itself swings. Asked why the mean-reversion speed is 2, she draws the
at-the-money term structure bending back to its long-run level within a year. Each of the
five numbers of the \omterm{def:dv:stochastic-volatility:heston}{Heston model} has a meaning on the desk, and that is why it has survived
thirty years of better-fitting alternatives. This chapter introduces stochastic volatility,
derives what can be derived about the \omterm{def:dv:stochastic-volatility:heston}{Heston model}, prices options with its characteristic
function, calibrates it to the surface of chapter 9, and reads its parameters the way a trader
does. It also measures where it fails: at short expiries its skew is half what the market
shows.

\section{Stochastic volatility models}

\omterm{def:dv:local-volatility:model}{Local volatility} makes volatility a function of the spot, and so ties every move of the smile to
a move of the spot. Real volatility has a life of its own: it rises after falls, but it also
rises and falls on days when the spot does not move. The leverage effect (One Quant Book 4,
chapter 18) is the first fact, measured as a negative correlation between returns and subsequent
volatility, strong and short-lived for indices and weaker for single stocks; volatility clustering
is the second.

\begin{definition}[Stochastic volatility model]\label{def:dv:stochastic-volatility:sv}
A \emph{stochastic volatility model}\index{stochastic volatility model} specifies, under the
pricing measure,
\[
dS_t=(r-q)S_t\,dt+\sqrt{v_t}\,S_t\,dW^1_t,\qquad dv_t=\alpha(v_t)\,dt+\beta(v_t)\,dW^2_t,\qquad
d\langle W^1,W^2\rangle_t=\rho\,dt,
\]
where the instantaneous variance $v_t$ is a second state variable. The
\emph{spot--volatility correlation}\index{spot--volatility correlation} $\rho$ ties the two
shocks together, and the scale of $\beta$ relative to $v$ is the \emph{volatility of
volatility}\index{volatility of volatility}.
\end{definition}

Two sources of randomness and one traded risky asset make the market incomplete (chapter 1): an
option cannot be replicated with the share and cash alone. It can with the share and a second
option, which hedges the variance risk. The drift $\alpha$ under the pricing measure absorbs a
market price of volatility risk that the real-world dynamics do not determine; calibration to
option prices, not a regression on history, fixes it.

\section{The Heston model}

\begin{definition}[Heston model]\label{def:dv:stochastic-volatility:heston}
The \emph{Heston model}\index{Heston model} takes the variance to be a square-root process (One Quant
Book 4, chapter 4):
\[
dv_t=\kappa(\bar v-v_t)\,dt+\eta\sqrt{v_t}\,dW^2_t,
\]
with five parameters: the initial variance $v_0$, the mean-reversion speed $\kappa$, the long-run
variance $\bar v$, the \omterm{def:dv:stochastic-volatility:sv}{volatility of volatility} $\eta$ and the correlation $\rho$. The variance stays
non-negative and never reaches zero when the Feller condition $2\kappa\bar v\ge\eta^2$ holds.
\end{definition}

The model's value is that it is affine: the log-price and the variance enter the pricing equation
linearly in the exponent, and its characteristic function is known in closed form.

\begin{proposition}[Heston characteristic function]\label{prop:dv:stochastic-volatility:cf}
Let $x_T=\ln(S_T/F_{0,T})$. Its characteristic function (One Quant Book 4, chapter 1) is
$\varphi_{x_T}(u)=\exp\bigl(C(u)+D(u)v_0\bigr)$, with $\xi=\kappa-\rho\eta iu$,
$d=\sqrt{\xi^2+\eta^2(iu+u^2)}$, $g=(\xi-d)/(\xi+d)$ and
\[
C=\frac{\kappa\bar v}{\eta^2}\Bigl((\xi-d)T-2\ln\frac{1-ge^{-dT}}{1-g}\Bigr),\qquad
D=\frac{\xi-d}{\eta^2}\,\frac{1-e^{-dT}}{1-ge^{-dT}} .
\]
\end{proposition}
\admitted

The functions $C$ and $D$ solve two ordinary differential equations (a Riccati equation and an
integral of it) obtained by substituting an exponential-affine guess into the pricing equation;
Heston solved them in 1993. The form above keeps the complex logarithm on its principal branch for
all $u$; Heston's original form, algebraically equal, jumps between branches for long expiries and
gives wrong prices, a defect known as the little Heston trap.

\section{The characteristic function and pricing}

With the characteristic function, a European option is one integral. In forward terms, with
$k=\ln(K/F)$,
\begin{equation}\label{eq:dv:stochastic-volatility:lewis}
C(K,T)=P(0,T)\Bigl(F-\frac{\sqrt{FK}}{\pi}\int_0^\infty\operatorname{Re}\bigl(e^{-iuk}
\varphi_{x_T}(u-\tfrac i2)\bigr)\frac{du}{u^2+\frac14}\Bigr),
\end{equation}
a formula that chapter 24 derives and generalises to any model with a known characteristic
function. All strikes of one expiry share the evaluation of $\varphi$, so a whole smile costs one
vectorised integral. Three checks guard the implementation: $\varphi(0)=\varphi(-i)=1$ (the
forward is a martingale); as $\eta\to0$ the smile flattens at the root-mean variance of the
deterministic path; and a simulation of the two equations reproduces the prices. On the chapter's
base parameters ($v_0=\bar v=0.04$, $\kappa=1.5$, $\eta=0.6$, $\rho=-0.7$) the one-year smile runs from
23.4\% at the 80 strike to 13.3\% at 120, and a 200\,000-path simulation agrees with the integral within
its noise.

\section{Calibration}

Calibrating Heston to a surface is a nonlinear least-squares problem in five parameters, solved by the
Levenberg--Marquardt algorithm (One Quant Book 4, chapter 24) on implied-volatility residuals, with
parameters transformed (logarithms, a hyperbolic tangent for $\rho$) so that no step leaves the
admissible set.

\begin{example}[Heston on the chapter 9 surface]\label{ex:dv:stochastic-volatility:calib}
Fitted to nine strikes (80 to 120) at four expiries (three months to two years) of the surface of
chapter 9, Heston returns $v_0=0.0267$ (a spot volatility of 16.3\%), $\kappa=2.05$, $\bar v=0.060$
(24.5\% long-run), $\eta=0.66$ and $\rho=-0.70$, with a root-mean-square error of 0.37 volatility point
and at most 0.7 point, at two years (\cref{fig:dv:stochastic-volatility:calib}). The Feller condition
fails, $2\kappa\bar v-\eta^2=-0.19$: the fitted variance can touch zero, as it does in most index
calibrations.
\end{example}

\begin{omfigure}
\begin{tikzpicture}
\begin{axis}[width=10.5cm,height=5cm,
  xlabel={strike},ylabel={implied volatility (\%)},
  tick label style={font=\small},label style={font=\small},
  xmin=80,xmax=120,ymin=12,ymax=28,grid=major,grid style={gray!25},
  legend columns=2,legend style={font=\footnotesize,at={(0.5,-0.3)},anchor=north,draw=none,fill=none,/tikz/every even column/.append style={column sep=8pt}},
  table/col sep=comma]
\addplot[omDef,very thick] table[x=strike,y=mkt3m]{figdata/derivatives/10-stochastic-volatility/calib.csv};
\addplot[omDef,thick,dashed] table[x=strike,y=mod3m]{figdata/derivatives/10-stochastic-volatility/calib.csv};
\addplot[omThm,very thick] table[x=strike,y=mkt2y]{figdata/derivatives/10-stochastic-volatility/calib.csv};
\addplot[omThm,thick,dashed] table[x=strike,y=mod2y]{figdata/derivatives/10-stochastic-volatility/calib.csv};
\legend{market 3 months,Heston 3 months,market 2 years,Heston 2 years}
\end{axis}
\end{tikzpicture}

\omcaption{The calibrated \omterm{def:dv:stochastic-volatility:heston}{Heston model} (dashed) against the surface it was fitted to (solid) at three
months and two years. Five parameters fit four expiries to within 0.7 volatility point. Data: the
tutorial.}\label{fig:dv:stochastic-volatility:calib}
\end{omfigure}

The fit hides a structural gap that shows at expiries outside the calibration set. Heston's at-the-money
skew tends to a finite limit as the expiry shrinks, while market skews keep steepening, roughly like
$T^{-1/2}$ in this surface. At one week the calibrated model's skew is half the market's
(\cref{fig:dv:stochastic-volatility:skew}); jumps (chapter 13) and rough volatility (chapter 12) are the
two answers.

\begin{omfigure}
\begin{tikzpicture}
\begin{loglogaxis}[width=10.5cm,height=5cm,
  xlabel={expiry (years, log scale)},ylabel={$-\partial\sigma_{\mathrm{imp}}/\partial k$ at the money},
  tick label style={font=\small},label style={font=\small},
  xmin=0.015,xmax=2.5,ymin=0.12,ymax=2,grid=major,grid style={gray!25},
  ytick={0.15,0.2,0.3,0.5,0.7,1,1.5},yticklabels={0.15,0.2,0.3,0.5,0.7,1,1.5},
  xtick={0.02,0.05,0.1,0.25,0.5,1,2},xticklabels={0.02,0.05,0.1,0.25,0.5,1,2},
  legend columns=2,legend style={font=\footnotesize,at={(0.5,-0.3)},anchor=north,draw=none,fill=none,/tikz/every even column/.append style={column sep=8pt}},
  table/col sep=comma]
\addplot[omThm,very thick,mark=*,mark size=1.5pt] table[x=t,y=market]{figdata/derivatives/10-stochastic-volatility/skew_term.csv};
\addplot[omDef,very thick,mark=square*,mark size=1.5pt] table[x=t,y=heston]{figdata/derivatives/10-stochastic-volatility/skew_term.csv};
\legend{market surface,calibrated Heston}
\end{loglogaxis}
\end{tikzpicture}

\omcaption{At-the-money skew by expiry, on log scales: the market surface's skew keeps steepening at short
expiries, the calibrated \omterm{def:dv:stochastic-volatility:heston}{Heston model}'s levels off near 0.7. At one week the model's skew is half the
market's; between three months and a year they agree, which is where the model was calibrated. Data: the
tutorial.}\label{fig:dv:stochastic-volatility:skew}
\end{omfigure}

\section{What the parameters mean to a trader}

Each parameter moves one feature of the surface
(\cref{fig:dv:stochastic-volatility:params,fig:dv:stochastic-volatility:term}):

\begin{itemize}
\item $\rho$ \emph{tilts} the smile: more negative, steeper skew. Moving $\rho$ from $-0.7$ to $-0.5$
lowers the one-year 90--110 skew from 5.71 to 3.94 volatility points.
\item $\eta$ \emph{bends} it: the wings rise with the \omterm{def:dv:stochastic-volatility:sv}{volatility of volatility}. The one-year curvature
measure $\tfrac12(\sigma_{90}+\sigma_{110})-\sigma_{100}$ is 0.30 point at $\eta=0.6$; a curvature of 0.5
point needs $\eta=0.77$.
\item $v_0$ and $\bar v$ set the short and long ends of the at-the-money term structure, and $\kappa$ the
speed at which one becomes the other; $\kappa$ also makes the skew fade with the expiry, since variance
shocks are forgotten at that rate.
\end{itemize}

\begin{omfigure}
\begin{tikzpicture}
\begin{groupplot}[group style={group size=2 by 1,horizontal sep=1.5cm},
  width=6.6cm,height=4.8cm,xlabel={strike},
  tick label style={font=\small},label style={font=\small},
  xmin=70,xmax=130,ymin=8,ymax=30,grid=major,grid style={gray!25},table/col sep=comma]
\nextgroupplot[ylabel={1-year implied volatility (\%)},title={\small correlation $\rho$},
  legend columns=3,legend style={font=\footnotesize,at={(0.5,-0.3)},anchor=north,draw=none,fill=none}]
\addplot[omThm,very thick] table[x=strike,y=rho_m09]{figdata/derivatives/10-stochastic-volatility/params.csv};
\addplot[omDef,thick] table[x=strike,y=rho_m05]{figdata/derivatives/10-stochastic-volatility/params.csv};
\addplot[omProp,thick,dashed] table[x=strike,y=rho_0]{figdata/derivatives/10-stochastic-volatility/params.csv};
\legend{$-0.9$,$-0.5$,$0$}
\nextgroupplot[title={\small volatility of volatility $\eta$},
  legend columns=3,legend style={font=\footnotesize,at={(0.5,-0.3)},anchor=north,draw=none,fill=none}]
\addplot[omProp,thick,dashed] table[x=strike,y=eta03]{figdata/derivatives/10-stochastic-volatility/params.csv};
\addplot[omDef,thick] table[x=strike,y=eta06]{figdata/derivatives/10-stochastic-volatility/params.csv};
\addplot[omThm,very thick] table[x=strike,y=eta09]{figdata/derivatives/10-stochastic-volatility/params.csv};
\legend{$0.3$,$0.6$,$0.9$}
\end{groupplot}
\end{tikzpicture}

\omcaption{One-year Heston smiles around the base parameters ($v_0=\bar v=0.04$, $\kappa=1.5$,
$\eta=0.6$, $\rho=-0.7$): the correlation rotates the smile (left), the \omterm{def:dv:stochastic-volatility:sv}{volatility of volatility} raises
its wings and, with a negative correlation, its skew (right). Data: the chapter's
code.}\label{fig:dv:stochastic-volatility:params}
\end{omfigure}

\begin{omfigure}
\begin{tikzpicture}
\begin{axis}[width=10.5cm,height=4.8cm,
  xlabel={expiry (years)},ylabel={at-the-money volatility (\%)},
  tick label style={font=\small},label style={font=\small},
  xmin=0,xmax=5,ymin=12,ymax=29,grid=major,grid style={gray!25},
  legend columns=3,legend style={font=\footnotesize,at={(0.5,-0.3)},anchor=north,draw=none,fill=none,/tikz/every even column/.append style={column sep=8pt}},
  table/col sep=comma]
\addplot[omDef,thick,mark=*,mark size=1.3pt] table[x=t,y=v0_002]{figdata/derivatives/10-stochastic-volatility/term.csv};
\addplot[gray,thick,mark=square*,mark size=1.3pt] table[x=t,y=v0_004]{figdata/derivatives/10-stochastic-volatility/term.csv};
\addplot[omThm,thick,mark=triangle*,mark size=1.6pt] table[x=t,y=v0_008]{figdata/derivatives/10-stochastic-volatility/term.csv};
\legend{$v_0=0.02$,$v_0=0.04$,$v_0=0.08$}
\end{axis}
\end{tikzpicture}

\omcaption{At-the-money term structures for three initial variances with the long-run variance at 0.04
(20\%) and $\kappa=1.5$: a calm market's upward-sloping term structure and a stressed market's inverted one,
both converging at the speed $\kappa$. With $v_0=\bar v$ the curve is not flat: it falls from 19.6\% to 16.9\% at
one year, since the at-the-money price is nearly linear in volatility, not in variance, and averaging over random
variance lowers the average volatility (the \omterm{prop:dv:stochastic-volatility:mixing}{mixing formula} below); the negative correlation lowers it further. Data:
the chapter's code.}\label{fig:dv:stochastic-volatility:term}
\end{omfigure}

For a trader the most useful identity is the one that holds without correlation.

\begin{proposition}[Mixing formula]\label{prop:dv:stochastic-volatility:mixing}
If $\rho=0$, a European option in a \omterm{def:dv:stochastic-volatility:sv}{stochastic volatility model} is worth the Black--Scholes price at the
path's root-mean variance, averaged over the variance paths:
$C=\E\bigl[C_{\mathrm{BS}}(\bar\sigma)\bigr]$ with $\bar\sigma^2=\frac1T\int_0^Tv_t\,dt$. With $\rho\ne0$ the same
holds after replacing the spot by an adjusted spot and the variance by $(1-\rho^2)$ times it; this is the
\emph{mixing formula}\index{mixing formula}.
\end{proposition}
\begin{proof}[Proof for $\rho=0$]
Conditionally on the whole variance path, $\ln S_T$ is normal with variance $\int_0^Tv_t\,dt$, since
$W^1$ is independent of $W^2$: the conditional price is Black--Scholes with that variance. Take the
expectation.
\end{proof}

The formula explains the smile's shape at once: Black--Scholes prices are convex in volatility away from
the money, so averaging over random volatility raises out-of-the-money prices more than at-the-money
ones, and the smile curves up (the \omterm{def:dv:greeks-and-the-hedging-pnl:greeks}{volga} of chapter 4 at work). A simulation of variance paths with
$\rho=0$ gives, through the formula, 7.136 for the one-year at-the-money call against 7.137 from the
characteristic function, well within its standard error of 0.011.

\section{Tutorial: pricing and calibrating Heston}\label{tut:dv:stochastic-volatility:tut}

\begin{tutorial}
\textbf{Goal.} Price with the characteristic function, check it three ways, calibrate to a surface and read
the parameters. \textbf{End state:} the four figures of the chapter and the numbers of
\cref{ex:dv:stochastic-volatility:calib}.
\begin{tutsteps}
\item \textbf{The characteristic function and the integral}, one smile at a time:
\omcode{code/firm/heston/firm_heston.py}{32}{54}{Heston characteristic function (principal-branch form) and the Fourier price of a smile.}\label{lst:dv:stochastic-volatility:cf}
\item \textbf{Levenberg--Marquardt} on implied-volatility residuals, with transformed parameters and a
damping that adapts to success:
\omcode{code/firm/heston/firm_heston.py}{74}{104}{Calibrating the five parameters.}\label{lst:dv:stochastic-volatility:calib}
\item \textbf{Run} \texttt{dv\_heston.calibration()}, \texttt{mixing\_check()} and \texttt{fig\_heston.py}.
\end{tutsteps}
\textbf{What to change next.} Replace the principal-branch form by Heston's original one and price a five-year
option with a large \omterm{def:dv:stochastic-volatility:sv}{volatility of volatility}: find the strikes where the price jumps; then add the one-week and
one-month expiries to the calibration and watch the fit degrade elsewhere.
\end{tutorial}

\section{Build: the Heston model}\label{bld:dv:stochastic-volatility:heston}

\begin{build}
\bfield{Purpose} The miniature firm's first model with its own volatility dynamics: it prices products that
depend on \omterm{def:dv:local-volatility:forward}{forward smiles} (chapter 16), it is the stochastic part of stochastic-local volatility (chapter 20),
and it is the test case of the Fourier and calibration engines (chapter 24).
\bfield{Interface} \texttt{Heston(v0, kappa, vbar, eta, rho)} with \texttt{cf(u, t)} and \texttt{feller()};
\texttt{call\_prices(model, fwd, strikes, t, df)}; \texttt{implied\_vols}; \texttt{calibrate(quotes, fwd,
start) -> (model, rmse)}; \texttt{simulate(model, s0, t, steps, n\_paths, seed)}.
\bfield{Rules} Characteristic function in the principal-branch form; the forward as numeraire; calibration on
implied-volatility residuals with transformed parameters; report the Feller gap, never impose it silently.
\bfield{Acceptance tests} \texttt{code/firm/heston/tests/}: $\varphi(0)=\varphi(-i)=1$; the small-$\eta$ limit
equals Black--Scholes at the deterministic root-mean variance; simulation agrees with the integral within
three standard errors; calibration recovers known parameters from their own smiles.
\bfield{Stretch} The quadratic-exponential simulation scheme (chapter 23) and analytic gradients of prices in the
parameters for a faster calibration.
\end{build}

\begin{omsources}
\item S.~L.~Heston, ``A closed-form solution for options with stochastic volatility with applications to bond
and currency options'', \emph{Review of Financial Studies} 6 (1993) 327--343.
\item H.~Albrecher, P.~Mayer, W.~Schoutens and J.~Tistaert, ``The little Heston trap'', \emph{Wilmott} (2007)
83--92.
\item J.~Hull and A.~White, ``The pricing of options on assets with stochastic volatilities'', \emph{Journal of
Finance} 42 (1987) 281--300.
\item M.~Romano and N.~Touzi, ``Contingent claims and market completeness in a stochastic volatility model'',
\emph{Mathematical Finance} 7 (1997) 399--412.
\item J.-P.~Bouchaud, A.~Matacz and M.~Potters, ``Leverage effect in financial markets: the retarded
volatility model'', \emph{Physical Review Letters} 87 (2001) 228701.
\item J.~Gatheral, \emph{The Volatility Surface}, Wiley, 2006, chapters 2--3.
\end{omsources}

\section{Exercises}

\begin{exercise}[$\star$]\label{exo:dv:stochastic-volatility:1}
Check the Feller condition for the base parameters and for the calibrated ones.
\end{exercise}

\begin{exercise}[$\star$]\label{exo:dv:stochastic-volatility:2}
With $v_0=0.08$, $\bar v=0.04$ and $\kappa=1.5$, give the deterministic variance after six months and one year,
and the at-the-money volatility of a one-year option as $\eta\to0$.
\end{exercise}

\begin{exercise}[$\star$]\label{exo:dv:stochastic-volatility:3}
Why must $\varphi_{x_T}(-i)=1$? What would a value of 1.01 mean?
\end{exercise}

\begin{exercise}[$\star\star$]\label{exo:dv:stochastic-volatility:4}
Why does the market price of volatility risk not appear in the calibrated model, and what does that imply about
using historical volatility to set $\bar v$?
\end{exercise}

\begin{exercise}[$\star\star$]\label{exo:dv:stochastic-volatility:5}
Use the \omterm{prop:dv:stochastic-volatility:mixing}{mixing formula} to explain why, with $\rho=0$, the Heston smile is symmetric in \omterm{def:dv:implied-volatility-and-its-surface:surface}{log-moneyness} around the
forward and why its at-the-money volatility lies below the root-mean variance.
\end{exercise}

\begin{exercise}[$\star\star$]\label{exo:dv:stochastic-volatility:6}
Which parameter would you move to fit a market in which the one-month skew steepened overnight while the one-year
skew did not? Why can Heston not fit it with one parameter move?
\end{exercise}

\begin{exercise}[$\star\star\star$]\label{exo:dv:stochastic-volatility:7}
\emph{Coding.} With \texttt{eta\_for\_butterfly}, find the \omterm{def:dv:stochastic-volatility:sv}{volatility of volatility} that gives a one-year curvature
of 0.5 volatility point, and the one-year 90--110 skew at that $\eta$.
\end{exercise}

\begin{exercise}[$\star\star\star$]\label{exo:dv:stochastic-volatility:8}
\emph{Find the flaw.} ``Our Heston calibration fits the three-month to two-year surface within 0.4 points, so we use
it to price weekly options on the index.''
\end{exercise}

\section{Problem: Reading the Parameters}

\begin{problem}[{Weekend problem --- what each Heston number does to the smile}]\label{pb:dv:stochastic-volatility:1}
A junior trader inherits the base \omterm{def:dv:stochastic-volatility:heston}{Heston model} ($v_0=\bar v=0.04$, $\kappa=1.5$, $\eta=0.6$, $\rho=-0.7$, zero rates)
and must adjust it to two changes in the market: the one-year 90--110 skew falls by about 1.8 points, and the one-year
curvature rises to 0.5 point.

\medskip\noindent\textbf{Part I --- The base smile.}
\begin{enumerate}
\item Give the one-year implied volatilities at the 80, 90, 100, 110 and 120 strikes.
\item Give the one-year 90--110 skew and the curvature measure.
\item Is the Feller condition satisfied?
\item Why is the at-the-money volatility below 20\% although $v_0=\bar v=0.04$?
\item What does the model predict for the at-the-money volatility at five years?
\end{enumerate}

\medskip\noindent\textbf{Part II --- The correlation.}
\begin{enumerate}[resume]
\item Give the 90--110 skew with $\rho=-0.5$.
\item Give the change in skew for the change of 0.2 in correlation.
\item Does the at-the-money volatility move when $\rho$ changes? Why?
\item Which products are most exposed to $\rho$?
\item What does $\rho=-0.7$ say about the index when it falls 1\% in a day?
\end{enumerate}

\medskip\noindent\textbf{Part III --- The \omterm{def:dv:stochastic-volatility:sv}{volatility of volatility}.}
\begin{enumerate}[resume]
\item Give the $\eta$ that produces a one-year curvature of 0.5 point.
\item What happens to the skew at that $\eta$, and why?
\item How would you restore the skew?
\item What is the Feller gap at that $\eta$?
\item Which products are most exposed to $\eta$?
\end{enumerate}

\medskip\noindent\textbf{Part IV --- Judgement.}
\begin{enumerate}[resume]
\item Why should the trader recalibrate all five parameters rather than move one at a time?
\item Why might the desk keep $\kappa$ fixed across days?
\item What does the model get wrong at one week?
\item State the \emph{named result}: the change in the one-year 90--110 skew for a change of $+0.2$ in the correlation,
and the \omterm{def:dv:stochastic-volatility:sv}{volatility of volatility} that matches a curvature of 0.5 point.
\item In one sentence: what does the \omterm{def:dv:stochastic-volatility:sv}{volatility of volatility} price?
\end{enumerate}
\end{problem}

\section{Interview questions}

\begin{interviewq}[$\star$ \iqroles{trader, researcher}]\label{iq:dv:stochastic-volatility:1}
Write the \omterm{def:dv:stochastic-volatility:heston}{Heston model} and say what each parameter does to the smile.
\end{interviewq}

\begin{interviewq}[$\star$ \iqroles{researcher}]\label{iq:dv:stochastic-volatility:2}
Why is a stochastic volatility market incomplete, and how do you hedge an option in it?
\end{interviewq}

\begin{interviewq}[$\star\star$ \iqroles{researcher}]\label{iq:dv:stochastic-volatility:3}
Derive the price of a European option in a \omterm{def:dv:stochastic-volatility:sv}{stochastic volatility model} with zero correlation.
\end{interviewq}

\begin{interviewq}[$\star\star$ \iqroles{developer, researcher}]\label{iq:dv:stochastic-volatility:4}
Your Heston pricer returns wrong prices for long expiries with a large \omterm{def:dv:stochastic-volatility:sv}{volatility of volatility}. What do you suspect?
\end{interviewq}

\begin{interviewq}[$\star\star$ \iqroles{trader, risk}]\label{iq:dv:stochastic-volatility:5}
Your calibrated Heston correlation moved from $-0.7$ to $-0.4$ overnight with the same fit error. What do you do?
\end{interviewq}

\begin{interviewq}[$\star\star\star$ \iqroles{researcher}]\label{iq:dv:stochastic-volatility:6}
Why is the Heston skew too flat at short expiries? Name two ways to fix it.
\end{interviewq}
